\subsection{Hyperparameter definitions and experimental scope}

Table~\ref{tab:hyperparameters} defines the LM hyperparameters reported in the
ablation tables. The parameters control three related stages: construction of
the local context, generation of candidate token partitions, and scoring or
regularization of the induced graph. Their interpretation is
method-dependent because not every LM implementation consumes every
parameter.

\begin{table*}[t]
  \centering
  \small
  \begin{tabular}{
    p{0.18\textwidth}
    p{0.12\textwidth}
    p{0.27\textwidth}
    p{0.28\textwidth}
  }
    \hline
    Parameter & Evaluated values & Operational definition &
    Interpretation \\
    \hline
    \texttt{lmWindowLength} ($W$)
    & $2,3,4$
    & Number of consecutive input tokens considered when training or
      selecting a locally coherent partition sequence.
    & Larger values provide broader contextual evidence but enlarge the
      candidate-combination space. The \texttt{lm-word} baseline retains its
      default value of $10$, although it performs no subword partitioning.
      \\

    \texttt{lmSlideLength} ($S$)
    & $2,3,4$
    & Number of partition units in the local sequence window accumulated by
      the transducer. Successive windows advance by one position.
    & Controls the order of the local sequential statistics; despite its
      name, it is a window length rather than the stride between windows.
      \\

    \texttt{lmTopSplit} ($K$)
    & $1,3,5$
    & Maximum number of candidate segmentations retained for each token
      before sentence-level combinations are scored.
    & Larger values increase candidate diversity and computational cost,
      whereas smaller values impose a more selective partition search.
      \\

    \texttt{lmSkip} ($D$)
    & $2,3,4$
    & Maximum backward distance used when adding or scoring skip
      dependencies between partition units.
    & Determines the contextual horizon of non-adjacent graph connections;
      larger values permit longer-range dependencies.
      \\

    \texttt{lmStemLength} ($L_{\min}$)
    & $7$
    & Minimum prefix length from which stem--remainder candidates are
      generated during dictionary construction.
    & Restricts very short stem candidates. Tokens no longer than this
      threshold remain available as unsplit whole-token candidates.
      \\

    \texttt{lmPrune} ($P$)
    & $100$
    & Maximum number of high-scoring outgoing transitions retained at each
      transducer or rank state when pruning is enabled.
    & Bounds graph branching and memory use while preserving the strongest
      observed transitions.
      \\

    \texttt{lmLikelihoodWeight} ($\lambda_{\mathrm{like}}$)
    & $0,.05,.15,$\newline $.30,.50$
    & Weight assigned to contextual transition likelihood in the
      \texttt{RankLM} score update.
    & Larger values increase the contribution of graph-based contextual
      evidence relative to the candidate's existing score.
      \\

    \texttt{lmPriorWeight} ($\lambda_{\mathrm{prior}}$)
    & $1,.95,.85,$\newline $.70,.50$
    & Weight assigned to the current or prior candidate score in
      \texttt{RankLM}.
    & Preserves evidence already associated with a candidate. In the
      evaluated grid,
      $\lambda_{\mathrm{prior}}=1-\lambda_{\mathrm{like}}$.
      \\

    \texttt{lmLengthPenalty}
    & none; inverse parts; inverse square-root parts; inverse log parts
    & Rescales a candidate score $s$ according to its number of non-empty
      partitions $m$.
    & Controls the influence of segmentation granularity on ranking, with
      progressively different normalization rates for multi-part candidates.
      \\
    \hline
  \end{tabular}
  \caption{Operational definitions of the LM hyperparameters used in the
  ablation experiments.}
  \label{tab:hyperparameters}
\end{table*}

For \texttt{RankLM}, the contextual update can be summarized as
\begin{equation}
  s_i^{\prime}
  =
  \lambda_{\mathrm{prior}}s_i
  +
  \lambda_{\mathrm{like}}\ell_i,
\end{equation}
where $s_i$ is the current candidate score and $\ell_i$ is its estimated
forward-transition likelihood. The tested length normalizations take the
form
\begin{equation}
  \widetilde{s}
  =
  \frac{s}{g(m)},
  \qquad
  g(m)\in
  \left\{
    1,\,
    m,\,
    \sqrt{m},\,
    \log(m+1)
  \right\}.
\end{equation}
The four choices correspond respectively to no penalty, inverse-parts,
inverse-square-root-parts, and inverse-log-parts normalization.

The supplementary \texttt{lm-subword} configurations also vary
\texttt{lmSample} over $5$, $10$, and $20$. This parameter limits the number
of entropy-ranked token partitions sampled after applying
\texttt{lmTopSplit}; its effective value is consequently bounded by the
number of retained candidates. The ranking iteration count is fixed at $10$
and is not part of the ablation grid.

Several design constraints limit one-parameter interpretation. In the base
grid, $W$, $S$, and $D$ are assigned the same value and therefore vary
together. Likewise, the likelihood and prior weights are complementary rather
than independently manipulated. Moreover, the likelihood/prior weighting and
length normalization are directly implemented by \texttt{RankLM};
\texttt{LemmaLM} records these settings for grid symmetry but does not
directly reference them in its current scoring path. Finally,
\texttt{lmStemLength} and \texttt{lmPrune} remain fixed at $7$ and $100$,
respectively, and their individual effects cannot be estimated from this
experiment.
